\documentclass[11pt]{article}
\usepackage[margin=1in]{geometry}
\usepackage{amsmath,amssymb}
\newcommand{\finalanswer}[1]{\par\medskip\noindent\fbox{\parbox{0.94\linewidth}{\textbf{Answer:} #1}}}
\title{STAT 412 Homework 6 --- Question 5}
\author{}\date{}
\begin{document}\maketitle
\section*{Problem}
For the 16 stock prices with known $\sigma=2.73$, construct 90\% and 99\% confidence intervals for $\mu$.
\section*{Work}
The sample mean is
\[
\bar x=\frac{\sum x_i}{16}=19.989375\approx19.99.
\]
Use $\bar x\pm z_{\alpha/2}\sigma/\sqrt n$.
For 90\%, $z=1.645$:
\[
E=1.645\frac{2.73}{\sqrt{16}}\approx1.12,
\quad
19.989375\pm1.1226=(18.87,21.11).
\]
For 99\%, $z=2.576$:
\[
E=2.576\frac{2.73}{4}\approx1.76,
\quad
19.989375\pm1.7580=(18.23,21.75).
\]
The correct interpretation is the confidence-interval interpretation: with 90\% confidence, the population mean closing stock price is between the 90\% interval endpoints; with 99\% confidence, it is between the 99\% interval endpoints. The 99\% confidence interval is wider because a higher confidence level uses a larger critical value.
\finalanswer{90\% CI $(18.87,21.11)$; 99\% CI $(18.23,21.75)$. Interpretation choice B. The 99\% interval is wider.}
\end{document}
